\section*{第 \ref{ch:g10:muscles-and-joints} 章——肌肉与关节}

\begin{solution}{exo:g10:muscles-and-joints:1}
骨端的\omterm{def:g10:muscles-and-joints:joint}{软骨}（使\omterm{def:g10:muscles-and-joints:joint}{关节}滑动平顺、分散载荷）；\omterm{def:g10:muscles-and-joints:joint}{关节}囊
（封闭\omterm{def:g10:muscles-and-joints:joint}{关节}）及其中的滑液（润滑）；\omterm{def:g10:muscles-and-joints:joint}{韧带}
（把骨连接在一起，限制运动）；膝\omterm{def:g10:muscles-and-joints:joint}{关节}中的半月板
（分散载荷）。
\end{solution}

\begin{solution}{exo:g10:muscles-and-joints:2}
\omterm{def:g10:muscles-and-joints:muscle}{肌腱}连接肌肉与骨，传递肌肉的拉力；\omterm{def:g10:muscles-and-joints:joint}{韧带}
连接\omterm{def:g10:muscles-and-joints:joint}{关节}两侧的骨，限制其运动。二者都
不会收缩。
\end{solution}

\begin{solution}{exo:g10:muscles-and-joints:3}
肌肉只能通过缩短产生拉力，不能推动。要使\omterm{def:g10:muscles-and-joints:joint}{关节}
向两个方向运动，需要屈肌和伸肌，二者各自
抵消对方所产生的作用。
\end{solution}

\begin{solution}{exo:g10:muscles-and-joints:4}
扭伤：\omterm{def:g10:muscles-and-joints:joint}{韧带}拉伸或撕裂。肌肉撕裂：\omterm{def:g10:muscles-and-joints:muscle}{肌纤维}
断裂。脱位：骨端被挤出\omterm{def:g10:muscles-and-joints:joint}{关节}，同时\omterm{def:g10:muscles-and-joints:joint}{关节}囊
和\omterm{def:g10:muscles-and-joints:joint}{韧带}拉伸或撕裂。
\end{solution}

\begin{solution}{exo:g10:muscles-and-joints:5}
慢\omterm{def:g10:muscles-and-joints:muscle}{肌纤维}：速度慢、力量弱、不易疲劳，\omterm{def:g10:cells-common-unit:organelle}{线粒体}和毛细血管丰富，
颜色发红，依靠呼吸作用供能。快\omterm{def:g10:muscles-and-joints:muscle}{肌纤维}：速度快、力量大，
约一分钟即疲劳，\omterm{def:g10:cells-common-unit:organelle}{线粒体}少，颜色发白，依靠磷酸盐储备
和\omterm{def:g10:cell-metabolism:fermentation}{发酵}供能。
\end{solution}

\begin{solution}{exo:g10:muscles-and-joints:6}
负荷 $2 \times 9.8 = \qty{19.6}{N}$；二头肌拉力 $19.6 \times 32/4
\approx \qty{157}{N}$。
\end{solution}

\begin{solution}{exo:g10:muscles-and-joints:7}
推力可能来自任一方向，而每块肌肉只能抵抗
一个方向的推力。两块肌肉同时收缩，使\omterm{def:g10:muscles-and-joints:joint}{关节}在两个方向上同时变硬；
无论推力来自哪个方向，其中一块肌肉已经在拉动，
与之相抗。
\end{solution}

\begin{solution}{exo:g10:muscles-and-joints:8}
\omterm{def:g10:muscles-and-joints:muscle}{肌纤维}缩短 \qty{0.6}{cm}；\omterm{def:g10:muscles-and-joints:muscle}{肌腱}使骨在
距肘 \qty{5}{cm} 处移动 \qty{0.6}{cm}，因此手移动
$0.6 \times 35/5 = \qty{4.2}{cm}$。
\end{solution}

\begin{solution}{exo:g10:muscles-and-joints:9}
活动范围大要求\omterm{def:g10:muscles-and-joints:joint}{关节}窝浅、\omterm{def:g10:muscles-and-joints:joint}{韧带}松弛；同样的松弛
使肱骨头在受力超出\omterm{def:g10:muscles-and-joints:joint}{韧带}所能维持的范围时
脱离\omterm{def:g10:muscles-and-joints:joint}{关节}窝。
\end{solution}

\begin{solution}{exo:g10:muscles-and-joints:10}
脱水及出汗造成的盐分流失，以及\omterm{def:g10:muscles-and-joints:muscle}{肌纤维}
在比赛末段的疲劳。处理方法：轻轻拉伸肌肉，并
饮用含盐的水。
\end{solution}

\begin{solution}{exo:g10:muscles-and-joints:11}
短跑运动员大腿中的快\omterm{def:g10:muscles-and-joints:muscle}{肌纤维}占 70\%，这一比例很大程度上
由遗传决定；训练可使纤维增粗、改善其功能配备，但
不会改变纤维的数量。未经训练的人快、慢\omterm{def:g10:muscles-and-joints:muscle}{肌纤维}
各半，其中慢\omterm{def:g10:muscles-and-joints:muscle}{肌纤维}还有更大的发展空间。
\end{solution}

\begin{solution}{exo:g10:muscles-and-joints:12}
$4 \times 70 \times 9.8 \approx \qty{2.7}{kN}$。应力 $2744/100
\approx \qty{27}{N/mm^2}$，约为破坏应力的四分之一，即
安全系数为 4。
\end{solution}

\begin{solution}{exo:g10:muscles-and-joints:13}
肌肉血供丰富，且含有可修复\omterm{def:g10:muscles-and-joints:muscle}{肌纤维}的储备\omterm{def:g10:cells-common-unit:cell}{细胞}，
修复只需数周。\omterm{def:g10:muscles-and-joints:joint}{韧带}血管少，营养物质和
修复\omterm{def:g10:cells-common-unit:cell}{细胞}到达缓慢，需要数月。\omterm{def:g10:muscles-and-joints:joint}{软骨}完全没有血管，
几乎没有可分裂的\omterm{def:g10:cells-common-unit:cell}{细胞}，仅靠滑液供给营养，
几乎无法修复。
\end{solution}

\begin{solution}{exo:g10:muscles-and-joints:14}
健美训练使纤维增粗（每根纤维中肌原纤维增多），而
拉伸则使肌肉–\omterm{def:g10:muscles-and-joints:muscle}{肌腱}单位变长，并提高\omterm{def:g10:muscles-and-joints:joint}{关节}囊的
耐受能力。纤维数量在生命早期即已固定；成年肌肉通过
改变现有纤维的大小而增长或缩小。
\end{solution}

\begin{solution}{exo:g10:muscles-and-joints:15}
上楼梯时，\omterm{def:g10:muscles-and-joints:joint}{关节}所承受的力为体重的数倍；较薄的
\omterm{def:g10:muscles-and-joints:joint}{软骨}将该压力传递给下方的骨，而骨中
有神经，因而会感到疼痛。\omterm{def:g10:muscles-and-joints:joint}{软骨}减少时，半月板
成为最后分散载荷的结构。膝周围强壮的肌肉可吸收部分
冲击并稳定\omterm{def:g10:muscles-and-joints:joint}{关节}，降低残存\omterm{def:g10:muscles-and-joints:joint}{软骨}承受的
峰值压力。
\end{solution}

\begin{solution}{pb:g10:muscles-and-joints:1}
\textbf{1.} $70 \times 9.8 = \qty{686}{N}$。

\textbf{2.} $686 \times 10/5 \approx \qty{1.4}{kN}$。

\textbf{3.} $686 \times 20/5 \approx \qty{2.7}{kN}$。

\textbf{4.} 分别为体重的两倍和四倍：每上一级台阶，股四头肌
所施加的拉力相当于四个人的重量，这也是长距离爬坡时大腿
感到酸痛的原因。

\textbf{5.} 股后肌群位于大腿后侧；它们
在股四头肌伸直膝\omterm{def:g10:muscles-and-joints:joint}{关节}时放松并被拉长，
同时轻微收缩以稳定\omterm{def:g10:muscles-and-joints:joint}{关节}。

\textbf{6.} 地面反作用力 $3 \times 686 \approx \qty{2.06}{kN}$；\omterm{def:g10:muscles-and-joints:muscle}{肌腱}
所受拉力 $2058 \times 20/5 \approx \qty{8.2}{kN}$，相当于
体重的十二倍。

\textbf{7.} $8232/120 \approx \qty{69}{N/mm^2}$，约为
断裂应力的三分之二。每次落地都耗用了\omterm{def:g10:muscles-and-joints:muscle}{肌腱}的大部分安全余量。

\textbf{8.} 载荷 $8.2 + 2.1 \approx \qty{10.3}{kN}$。有半月板时：
$10300/600 \approx \qty{17}{N/mm^2}$；无半月板时：$10300/100 \approx
\qty{103}{N/mm^2}$，约为前者的六倍。

\textbf{9.} 每次落地时，同一\omterm{def:g10:muscles-and-joints:joint}{软骨}所承受的压强为原来的六倍，
表面被压溃的速度超过了其修复速度，于是
逐渐磨损。

\textbf{10.} \omterm{def:g10:muscles-and-joints:joint}{软骨}之间有一层液体薄膜，
使两者在弯曲时可以无干摩擦地滑动，
并对冲击起轻微的缓冲作用。

\textbf{11.} 胫骨在股骨下方发生扭转，并伴有侧向
弯曲：这是膝\omterm{def:g10:muscles-and-joints:joint}{关节}结构无法承受的转动，\omterm{def:g10:muscles-and-joints:joint}{韧带}
的作用正是抵抗这种运动。

\textbf{12.} 股四头肌只沿小腿方向施力，使膝\omterm{def:g10:muscles-and-joints:joint}{关节}伸直；
它对扭转并无拮抗作用。而且，损伤在
几百分之一秒内即已发生，快于任何肌肉反应。

\textbf{13.} 血液来自撕裂\omterm{def:g10:muscles-and-joints:joint}{韧带}和\omterm{def:g10:muscles-and-joints:joint}{关节}囊的血管，
另加额外的滑液：\omterm{def:g10:muscles-and-joints:joint}{关节}囊被刺破或发炎，
\omterm{def:g10:muscles-and-joints:joint}{关节}因此积液。

\textbf{14.} \omterm{def:g10:muscles-and-joints:joint}{韧带}血供不足，需要数月
时间，且往往需要手术；撕裂的半月板是几乎没有血管的\omterm{def:g10:muscles-and-joints:joint}{软骨}，
通常无法自行愈合，撕裂部分需
修剪或缝合。

\textbf{15.} 股后肌群是股四头肌的拮抗肌，牵拉
胫骨向后；两者同时强力收缩，可使膝\omterm{def:g10:muscles-and-joints:joint}{关节}变得
坚硬，抵抗撕裂\omterm{def:g10:muscles-and-joints:joint}{韧带}已无法制约的胫骨前移，
在手术前后保护\omterm{def:g10:muscles-and-joints:joint}{关节}。

\textbf{16.} 每天约 4000 次，每次受力 \qty{1.4}{kN} 或更大。

\textbf{17.} 每年约 $\num{1.5e6}$ 次负荷。\omterm{def:g10:muscles-and-joints:muscle}{肌腱}炎是损伤
累积的速度超过了\omterm{def:g10:muscles-and-joints:muscle}{肌腱}组织缓慢修复的速度；只有休息
才能让修复赶上，因为力量增强并不能降低
每一步的负荷。

\textbf{18.} 落地共计 $40 \times 8.2 \approx \qty{330}{kN}$；
台阶为 $4000 \times 1.4 \approx \qty{5600}{kN}$——台阶承受的总负荷
约为落地的十七倍，但真正威胁断裂的
是落地时的峰值负荷。

\textbf{19.} 地面反作用力加倍，每一步中的\omterm{def:g10:muscles-and-joints:muscle}{肌腱}和\omterm{def:g10:muscles-and-joints:joint}{关节}
受力也随之加倍；骑自行车和游泳锻炼肌肉，
同时使膝\omterm{def:g10:muscles-and-joints:joint}{关节}所受负荷远低于体重。

\textbf{20.} 股四头肌\omterm{def:g10:muscles-and-joints:muscle}{肌腱}在上台阶时承受体重的四倍，
落地时约为十二倍；\omterm{def:g10:muscles-and-joints:joint}{韧带}从未被这些负荷考验，
却正是扭转动作暴露出的薄弱组织。
\end{solution}
